\section{Discusión Crítica de Hallazgos}
\label{sec:discusion}

\subsection{Teoría asintótica y tiempo de pared}
La teoría ofrece una jerarquía clara. Para el Fibonacci heap,
\begin{equation}
    T_{\text{Fib}}=O(|E|+|V|\log|V|),
    \label{eq:fib}
\end{equation}
mientras que para un heap binario indexado,
\begin{equation}
    T_{\text{Bin}}=O((|V|+|E|)\log|V|).
    \label{eq:bin}
\end{equation}
La primera cota es menor cuando $|E|$ domina a $|V|$. Sin embargo, la notación $O(\cdot)$ oculta las constantes y el costo de las operaciones que no son comparaciones. En la implementación evaluada, el Fibonacci heap debe crear objetos, enlazar listas circulares, cortar nodos y consolidar raíces. Esas operaciones tienen un costo real que no aparece en~\eqref{eq:fib}. El resultado experimental —Fibonacci nunca gana y es hasta 1.59 veces más lento— no refuta la teoría: muestra que la cota asintótica describe el crecimiento, no el tiempo en una máquina y un lenguaje concretos.

La evidencia es consistente con \textcite{cherkassky1996shortest} y \textcite{lamarca1996cache}: en redes dispersas, los heaps contiguos suelen superar a las estructuras enlazadas porque sus ventajas de localidad y de asignación pesan más que la mejora amortizada de \texttt{decrease\_key}. El Fibonacci heap sería competitivo si el número de decrementos fuera enorme frente al número de extracciones y si la asignación de nodos fuera barata. En las instancias medidas, los decrementos efectivos fueron 6695 frente a 5000 extracciones en $|V|=5000$ disperso, y solo 2976 frente a 800 extracciones en $|V|=800$ denso. La relación no es la de un régimen donde el término $|E|$ de~\eqref{eq:fib} domine con holgura.

\subsection{Por qué el heap 4-ario gana en algunas escalas}
El heap 4-ario reduce la altura de $\lceil\log_2|V|\rceil$ a $\lceil\log_4|V|\rceil$. En $|V|=5000$, la altura pasa de 13 a 7 niveles. Además, los cuatro hijos de un nodo ocupan posiciones contiguas y suelen caber en una misma línea de caché. Esa combinación explica por qué el 4-ario fue hasta 24.4\,\% más rápido en $|V|=500$ y 16.5\,\% en $|V|=5000$. Pero la ventaja no es gratuita: \texttt{sift-down} compara cuatro hijos por nivel y, en escalas pequeñas, el costo de esas comparaciones y de la aritmética de índices puede superar el ahorro de altura. Por eso el 4-ario fue ligeramente más lento en $|V|=100$ y $|V|=2500$ del perfil disperso, y en $|V|=800$ del denso.

El hecho de que los IQR se solapen en las 10 escalas impide afirmar que el 4-ario sea mejor incluso donde su mediana es menor. El solapamiento no significa que las colas sean idénticas; significa que la variabilidad entre ejecuciones es del mismo orden que la diferencia entre medianas. Para declarar una ventaja estable harían falta más repeticiones, aislamiento de CPU o una carga de trabajo más larga. Con el diseño actual, la conclusión correcta es condicional: el 4-ario es un candidato, no un ganador universal.

\subsection{Por qué el heap binario gana en la consulta OSM}
En la consulta OSM, el binario fue 26.9\,\% más rápido que el 4-ario y 66.7\,\% más rápido que el Fibonacci. La red tiene solo 891 vértices y 2139 aristas, y la consulta punto a punto termina antes de explorar toda la componente. En ese régimen, el número de operaciones es pequeño y las constantes de implementación pesan más que la altura del árbol. El binario tiene la implementación más simple, menos comparaciones por nivel y menos aritmética de índices que el 4-ario; además, su IQR es amplio, lo que sugiere que el tiempo está dominado por el ruido del sistema y por el costo de relajación, no por la cola.

La comparación OSM no contradice el benchmark sintético: lo complementa. El 4-ario gana en escalas medianas y grandes del banco, pero pierde en una consulta pequeña y concreta. Esto refuerza la tesis de que la elección depende del tamaño, la densidad y el entorno. También advierte sobre un error frecuente: extrapolar el resultado de un benchmark sintético a una consulta real sin medirla.

\subsection{El papel de la densidad}
La densidad aumenta el número de intentos de relajación, pero no necesariamente el número de \texttt{decrease\_key}. En el perfil denso con $|V|=800$, hubo 25527 intentos y solo 2976 decrementos (11.7\,\%). La mayoría de los intentos no mejora la distancia tentativa y se descarta con una comparación. Por eso el Fibonacci heap, que optimiza \texttt{decrease\_key}, tiene poco margen de mejora. El 4-ario, que optimiza la extracción y la localidad, sí puede aprovechar la mayor altura relativa del árbol. Aun así, la diferencia entre 4-ario y binario en el perfil denso fue solo del 4.6\,\% en la mediana, con IQR solapados.

\subsection{Extensión $A^*$}
$A^*$ redujo las expansiones de 887 a 823 (7.2\,\%) sin alterar el costo. La reducción es menor que la que se obtiene cuando la heurística es muy informativa, porque $\lambda\lVert x_v-x_t\rVert$ es una cota inferior geométrica y no incorpora la estructura de la red ni los sentidos únicos. La heurística es consistente, y eso garantiza que no se reabran vértices cerrados; la optimalidad se mantiene. El resultado es útil como extensión, pero no convierte a $A^*$ en el objeto principal del estudio: la comparación de colas se hizo con Dijkstra para aislar el efecto de la estructura.

\subsection{Sensibilidad climática}
El factor climático uniforme no puede cambiar la ruta. Si $w_m(e)=m\,w(e)$ para todo $e$ y $m>0$, entonces para dos caminos $P$ y $Q$,
\begin{equation}
    w_m(P)<w_m(Q)\iff m\,w(P)<m\,w(Q)\iff w(P)<w(Q).
    \label{eq:invariancia}
\end{equation}
La ruta óptima es la misma; solo cambia el costo total. El incremento del 6.91\,\% es, por tanto, una medida de sensibilidad del tiempo estimado, no un cambio de decisión. Para modelar un efecto real del clima sobre la ruta, el factor tendría que depender del tramo: por ejemplo, penalizar más las calles con pendiente, las vías sin pavimento o las expuestas al viento. Esa extensión queda fuera del alcance de este proyecto y se propone como trabajo futuro.

\subsection{Marco de decisión}
La evidencia permite formular una regla práctica:
\begin{enumerate}
    \item \textbf{Valor por defecto: heap binario.} Es simple, robusto y ganó la consulta OSM y 3 de 10 escalas sintéticas. Para una aplicación nueva, es la opción con menor riesgo de implementación.
    \item \textbf{Candidato: heap 4-ario.} Si un benchmark local, con la distribución real de consultas, muestra una ventaja consistente y los IQR no se solapan, el 4-ario puede adoptarse. No se justifica por la teoría sola.
    \item \textbf{Fibonacci: uso teórico o experimental.} Su cota es valiosa para analizar el peor caso y para escenarios con muchísimos decrementos y asignación barata. En Python, no fue competitivo.
    \item \textbf{$A^*$: consultas punto a punto con coordenadas.} Reduce expansiones sin perder optimalidad; su magnitud depende de la fuerza de la heurística.
\end{enumerate}

\subsection{Limitaciones}
\begin{itemize}
    \item \textbf{Una sola máquina y un solo lenguaje.} Los resultados corresponden a Python 3.12 en un Ryzen 5 5600H. En C++ o Rust, las constantes y posiblemente el orden cambiarían.
    \item \textbf{Sin métricas de memoria.} No se midió RSS ni asignaciones; las comparaciones espaciales son teóricas.
    \item \textbf{Sin pruebas de significancia.} Con distribuciones no normales y 45 muestras por celda, se usaron medianas e IQR. No se reportan valores $p$ ni intervalos de confianza.
    \item \textbf{Instancias sintéticas.} El generador produce grafos aleatorios alcanzables, no redes viales geométricas. La red OSM mitiga parcialmente esta limitación, pero solo con una consulta.
    \item \textbf{Clima de modelo.} Open-Meteo entrega datos de modelo, no de una estación física; el factor es un escenario uniforme, no una calibración de tránsito.
    \item \textbf{Instancia ilustrativa.} La red de seis nodos tiene nombres reales, pero longitudes, pendientes y topología esquemáticas; no es un levantamiento geográfico.
\end{itemize}

\subsection{Validez externa}
La conclusión más robusta no es un ranking fijo, sino una metodología: medir la cola en la carga de trabajo objetivo, reportar dispersión y no extrapolar de un perfil a otro. El resultado puede cambiar con la implementación, la CPU o la distribución de consultas. Por eso el sistema deja el benchmark y los datos versionados: la decisión debe repetirse en el entorno de despliegue, no heredarse de este informe.
